\section{Motivations}

It is the business logic and the user that defines our software product to the outside world.
When considering quality requirements, products that have functionally have nothing in common are facing the exact same challenges.
This means they can benefit from the same solutions.
While companies may not have anything in common externally, they can benefit from the same scalability and extensibility patterns internally.
All companies are expected to be highly available and reliable.

\section{Cloud computing}

Provides infinite computation, storage and network capacity.
Infrastructure as a service, without this every company would spend months hiring specialists and building infrastructures to deliver a basic product.
Cloud charges only for the services that are used.
This removes the barriers for startups since there is no high upfront cost.
Cloud also provides message brokers, DBs etc.

\subsection{Downside}

None of that infrastructure is owned by the company.

It is out of the companies control.

Typically, they don't get the best hardware and a lot of it has already been used.
Runs the risk of failing and there is nothing that can be done about it.
The bigger the system the higher the probability of failure is.
Monthly bill will get bigger and bigger as the company grows.
As an architect we need to manage cost, scalability and profitability.

Cloud architecture can be thought of as building reliable systems using unreliable components.

This type of architecture will have to account for reliability and handling failures a lot more seriously.

\section{Scalabilty patterns}

\subsection{Motivations}
These patterns allow us to architect important patterns that allow is to architect and deploy a highly scalable system for every domain and use case.
This will enable out business to handle billions of requests per day process, petabytes of data and most importantly and do it in a cost-effective way in a cloud environment.